\documentclass{article}
\usepackage[utf8]{inputenc}
\usepackage{graphicx}
\usepackage{amsmath}
\usepackage{xcolor}
\usepackage{nopageno}
\usepackage{hyperref}

\newcommand{\pd}[3]{\frac{\partial^{#3} {#1}}{\partial {#2}^{#3}}}
\newcommand{\fpd}[2]{\frac{\partial {#1}}{\partial {#2}}}
\newcommand{\diff}[2]{\frac{d {#1}}{d {#2}}}


\begin{document}
\noindent You and a friend both measure something (for example how long you wait in line, how many cars pass a given intersection in a 20 minute period, how long a nail is, the amount of soda in a bottle, etc.). 

Do you think you will get the same value (measurements) every time?

Of course, not. Measurements have variation. This variation can sometimes be explained by error in the way the measurement is taken (measurement error). But even if we were able to measure perfectly, some measurements still have variation. There are some common models of measurements. Here are a few, each with an accompanying story. A random variable can be described by these or other stories.
\begin{itemize}
\item Binomial Distribution: This is a model for the number of ``successes" in $n$ independent Bernoulli trials (success/failure experiment) each with probability $p$ of ``success".\\
Example 1:\\
Example 2:\\

\item HyperGeometric Distribution: This is a model for the number of ``successes" in $n$ Bernoulli trials (success/failure experiment) from a finite population, so the trial are not independent, each trial has probability $p$ of ``success". \\
Example 1:\\
Example 2:\\

\item Poisson Distribution: This is a model for the number of (independent) events that occur in a fixed time (or space) interval, given the events occur at a constant ``probabilistic" rate.\\
Example 1:\\
Example 2:\\
\item Uniform Distribution: This is a model where each value in an interval is equally likely to be the measure. 
Example 1:\\
Example 2:\\

\item Exponential Distribution: This is a model of the amount of time (or space) until an event occurs, given the events occur at a constant ``probabilistic" rate. This is a model that assumes a lack of memory about how long you have already waited which means the number of events in disjoint time intervals must be independent.\\
Example 1:\\
Example 2:\\

\item Normal Distribution: \\
Example 1:\\
Example 2:\\
\end{itemize}

%Here are the formulas for the probability mass (if discrete random variable) or probability density (if continuous random variable) functions.
Here are the formulas for the probability mass (if random variable is discrete) or probability density (if random variable is continuous) functions.
\begin{itemize}
\item Binomial Distribution: two parameters $n$ and $p$\\
$$
p(x) = Pr(X = x) = \left\{
        \begin{array}{ll}
            \frac{n!}{x!(n-x)!}p^x(1-p)^{(n-x)} & \quad x = 0, 1, 2, 3, ... , n \\
            0 & \quad \text{otherwise}
        \end{array}
    \right.
$$
\item HyperGeometric Distribution: three parameters $N$, $R$, and $n$\\
$$
p(x) = Pr(X = x) = \left\{
        \begin{array}{ll}
            %\frac{\dbinom{R}{x} \dbinom{N-R}{n-x}}{\dbinom{N}{n}}& \quad \max(0,R+n-N) \leq x \leq \min(n,R) \\
            \frac{\left(\frac{R!}{x!(R-x)!}\right) \left(\frac{(N-R)!}{(n-x)!(N-R-n+x)!}\right)}{\left(\frac{N!}{n!(N-n)!}\right)}& \quad \max(0,R+n-N) \leq x \leq \min(n,R) \\
            0 & \quad \text{otherwise}
        \end{array}
    \right.
$$
\item Poisson Distribution: one parameter $\lambda$\\
$$
p(x) = Pr(X = x) = \left\{
        \begin{array}{ll}
            e^{-\lambda}\frac{\lambda^x}{x!} & \quad x = 0, 1, 2, 3, ... , n \\
            0 & \quad \text{otherwise}
        \end{array}
    \right.
$$
\item Uniform Distribution: two parameters $a$ and $b$\\
$$
f(x) = \left\{
        \begin{array}{ll}
            \frac{1}{b-a} & \quad a < x < b \\
            0 & \quad \text{otherwise}
        \end{array}
    \right.
$$
\item Exponential Distribution: one parameter $\lambda$\\
$$
f(x) = \left\{
        \begin{array}{ll}
            \lambda e^{-\lambda x} & \quad x > 0 \\
            0 & \quad x \leq 0
        \end{array}
    \right.
$$
\item Normal Distribution: two parameters $\mu$ and $\sigma^2$\\
$$
f(x) = \frac{1}{\sqrt{2\pi\sigma^2}}e^{-\frac{(x-\mu)^2}{2\sigma^2}}
$$
\end{itemize}

\noindent  Just like we did with more familiar models, we can explore these new models by changing parameter values and graphing these function in R. For these functions we can write our own function (user built function) just like we did before. We can also use the preprogramed functions that come as part of R. 

One of the course objectives for Math 119 is to read, write, and manipulate functions for use in data modeling. This assignment/activity is an opportunity to explore some common functions used to model measurements.


\end{document}